% Brief (grammar: small state machine with guarded exits, badges keyed to a
%   legend; atlas V/fig:anchor-card-lifecycle prescribes a state machine, not
%   a timeline -- the current fragment and the tikz-diagram-craft example
%   draft (redraw-anchor-card-lifecycle.tex) both draw a timeline of
%   admissible intervals instead, and the draft also names the wrong hue
%   (pdindigo; this chapter is pdcobalt). Redrawn to the atlas's own grammar
%   here; flagged for the author to reconcile the atlas row or the draft.
% Reader question: what are the only valid exits from an active card?
% Claim: an issued card is absorbed into exactly one of two terminal states,
%   revoked or expired; neither has an outgoing edge, so recovery is a new
%   card, never a transition back to active.
% Evidence: def:revocation, prop:no-reanim ("a revoked card cannot be
%   un-revoked within its TTL... they must issue a new card"),
%   prop:filter-mono (expiry only after t0+tau) [internal].
% Must distinguish: the two absorbing terminal states (revoked, expired) from
%   the one non-transition available after either -- issuing a new card,
%   drawn as a correction off-ramp, never a return arrow to issued.
% Grammar chosen: compact UML state machine, with the same-card lifecycle in
%   one panel and fresh issuance outside it; guards are keyed to a ruled
%   legend. Rejected: a timeline, which answers when rather than which exits,
%   and a tall fork, which spends vertical space without adding state meaning.
% Five-second test: "issued has exactly two guarded exits, both absorbing;
%   the only way out of either is a new card, not a returning arrow."
\begin{figure}[H]
\centering
\pdfigurehue{pdcobalt}%
\begin{tikzpicture}[pd figure,x=1cm,y=1cm,
  terminal/.style={pd breach state,double=pdpage,double distance=1.1pt,
    line width=.5pt,rounded corners=4.5pt}]
  \node[pd uml initial] (start) at (.20,2.00) {};
  \node[pd uml focus state,minimum width=2.05cm] (issued) at (1.75,2.00)
    {issued};
  \node[pd uml choice] (fork) at (3.65,2.00) {};
  \node[terminal,minimum width=2.25cm] (revoked) at (5.70,3.00)
    {revoked};
  \node[terminal,minimum width=2.25cm] (expired) at (5.70,1.00)
    {expired};

  \node[pd panel,fit=(start)(issued)(fork)(revoked)(expired),
        inner sep=6pt] (samecard) {};
  \node[pd kind,anchor=south] at (3.10,3.72) {same card ID};
  \draw[pd uml transition,shorten >=3pt] (start) -- (issued);
  \draw[pd uml transition,shorten >=3pt] (issued) -- (fork);
  \node[pd badge] (g1) at (4.60,2.62) {1};
  \node[pd badge] (g2) at (4.60,1.38) {2};
  \draw[pd breach rule,pd to badge] (fork) -- (g1);
  \draw[pd uml breach transition,pd from badge,shorten >=3pt] (g1) -- (revoked);
  \draw[pd breach rule,pd to badge] (fork) -- (g2);
  \draw[pd uml breach transition,pd from badge,shorten >=3pt] (g2) -- (expired);

  \node[pd ready state,minimum width=2.35cm] (fresh) at (9.15,3.00)
    {new card\\new ID};
  \draw[pd ready arrow,dash pattern=on 3.2pt off 2.2pt,shorten <=3pt,shorten >=3pt]
    (samecard.east |- revoked.east) -- (fresh.west);
  \node[pd note,anchor=north west,align=left,text width=3.2cm] at (7.22,2.45)
    {not a transition of the revoked card};
\end{tikzpicture}
\vspace{2pt}
{\pdfiglabel
\begin{tabular}{@{}c p{.82\linewidth}@{}}
\toprule
{\pdfiglabelbold Guard} & {\pdfiglabelbold Predicate} \\
\midrule
\raisebox{-.5ex}{\tikz[baseline=-.6ex]\node[pd badge]{1};} & explicit revocation: operator action or detected compromise \\
\raisebox{-.5ex}{\tikz[baseline=-.6ex]\node[pd badge]{2};} & natural expiry: issue time plus TTL \\
\bottomrule
\end{tabular}}
\caption{A UML state machine gives one card ID exactly two guarded exits,
both terminal. Recovery starts a new lifecycle on a fresh ID; it never
reanimates the revoked card. [internal]}
\label{fig:anchor-card-lifecycle}
\end{figure}
